\chapter{Getting started with AEDL}\label{ch:aedl-intro}

\section{The concept}

An experiment in GLoBES has one neutrino source (a point source, stationary
in time), one baseline with its matter density profile, and one detector. The
detector is described by its response to neutrinos of each flavour: cross
sections, energy resolution, efficiencies and backgrounds. The building blocks
are
\begin{description}
\item[fluxes] (\code{nuflux}) of the source, from a table or built in;
\item[cross sections] (\code{cross}), from tables;
\item[energy resolutions] (\code{energy}), a resolution function or a matrix;
\item[channels] (\code{channel}): one flux, one initial and one final flavour,
  one cross section and one energy resolution, with efficiencies;
\item[rules] (\code{rule}): signal and background, each a weighted sum of
  channels, with a $\chi^2$ function for systematics.
\end{description}
The $\chi^2$ of an experiment is the sum of the $\chi^2$ of its rules, and the
$\chi^2$ of a setup the sum over its experiments (Figure~\ref{fig:concept}).
Chapter~\ref{ch:rates} gives the formulas.

\begin{figure}[t]
\centering
\begin{tikzpicture}[>={Stealth[length=2mm]},font=\small,
  ch/.style={draw,rounded corners=2pt,fill=green!8,minimum width=17mm,minimum height=7mm},
  ru/.style={draw,rounded corners=2pt,fill=orange!12,minimum width=24mm,minimum height=8mm,align=center},
  ex/.style={draw,rounded corners=2pt,fill=blue!8,minimum width=40mm,minimum height=8mm}]
\node[ch] (c1) at (0,0) {channel 1};
\node[ch] (c2) at (2.2,0) {channel 2};
\node[ch] (c3) at (4.4,0) {channel 3};
\node[ch] (c4) at (6.6,0) {channel 4};
\node[ch] (c5) at (8.8,0) {channel 5};
\node[ru] (r1) at (1.1,-1.6) {rule 1\\$s$, $b$, $\chi^2_1$};
\node[ru] (r2) at (5.5,-1.6) {rule 2\\$s$, $b$, $\chi^2_2$};
\node[ru] (r3) at (8.8,-1.6) {rule 3\\$s$, $b$, $\chi^2_3$};
\node[ex] (e) at (5,-3.2) {experiment: $\chi^2=\chi^2_1+\chi^2_2+\chi^2_3$};
\draw[->] (c1) -- (r1); \draw[->] (c2) -- (r1); \draw[->] (c3) -- (r2);
\draw[->] (c4) -- (r2); \draw[->] (c2) -- (r2); \draw[->] (c5) -- (r3); \draw[->] (c4) -- (r3);
\draw[->] (r1) -- (e); \draw[->] (r2) -- (e); \draw[->] (r3) -- (e);
\end{tikzpicture}
\caption{Channels, rules and experiments. A channel can be used in several
rules, as signal in one and as background in another.}
\label{fig:concept}
\end{figure}

\section{A simple example}

The file \file{test/data/Simple.glb} describes a superbeam with a water
Cherenkov detector at 295~km:

\VerbatimInput[breaklines,frame=leftline,framerule=1.5pt,rulecolor=\color{orange!70!black},
  xleftmargin=6pt,framesep=8pt,fontsize=\small,numbers=left,numbersep=4pt]{../../test/data/Simple.glb}

\noindent It defines a flux from a table (lines 9--14), the detector mass, the
energy bins and sampling points (lines 17--26), the baseline and matter
profile (lines 29--30), two cross sections, a Gaussian energy resolution, four
channels, and two rules. The channel \code{\#mu\_e\_eff} has an energy dependent
efficiency computed in AEDL by interpolation (lines 47--53). Loaded with
\begin{julia}
simple = Experiment("Simple.glb")
\end{julia}
it can be used like any other experiment.

\section{Syntax}\label{sec:aedl-syntax}

\paragraph{File structure.} An AEDL file starts with the eight characters
\code{\%!GLoBES}. It should declare the GLoBES version it is written for,
\code{\$version = "3.0.0"}. Comments are as in C, \code{/* ... */} and
\code{// ...} to the end of the line.

\paragraph{Global and local variables.} Variables starting with \code{\$} are
global properties of the experiment (\code{\$baseline = 295}). Variables
starting with \code{@} belong to an environment
(\code{@time = 5} inside \code{nuflux}). Numbers are converted to integers
where integers are expected. Lists are written in braces,
\code{\$densitytab = \{2.8, 3.5\}}.

\paragraph{Environments and names.} Fluxes, cross sections, energy
resolutions, channels and rules are environments,
\begin{aedl}
channel(#name)<
   @channel = ...
>
\end{aedl}
named with a \code{\#}. Names are numbered in the order of their first
appearance within each type (starting at 1 in GLoBES.jl), and must be unique
across all types. A name can be declared before it is defined with an empty
environment, \code{channel(\#later)<>}. \code{rule\_names} and
\code{channel\_names} return the names of an experiment.

\paragraph{User variables and expressions.}\label{sec:aedl-variables}
Variables without prefix are constants for use in expressions:
\begin{aedl}
L = 1300
$baseline = L
$emax = 2 * sqrt(L / 100) + 0.5^2
\end{aedl}
The operators are \code{+ - * / \^}, and the functions \code{sin}, \code{cos},
\code{tan}, \code{asin}, \code{acos}, \code{atan}, \code{log}, \code{log10},
\code{exp} and \code{sqrt}, which apply to every element of a list. Variables
can also be defined by the program: \code{Experiment(file; L=810)}
(Chapter~\ref{ch:runtime}). Lists are stored in variables starting with
\code{\%}, \code{\%effs = \{0.2, 0.4, 0.6\}}.

\paragraph{List functions.}
\begin{description}[style=nextline,leftmargin=3cm]
\item[\code{bincenter()}] the centres of the energy bins defined so far;
\item[\code{samplingbincenter()}] the centres of the sampling points;
\item[\code{interpolation(x, y, k, z)}] the values at \code{z} of the linear
  ($k=1$) or natural cubic spline ($k=2$) interpolation of the points
  $(x, y)$;
\item[\code{copy(l)}] a copy of a list, needed when a list is used twice;
\item[\code{reverse(l)}] the list in reverse order;
\item[\code{echo(x)}, \code{echon(x)}, \code{line(n)}] print values (for
  debugging AEDL files), without or with line feeds, or $n$ line feeds.
\end{description}

\paragraph{Includes and several experiments.} \code{include "file.glb"}
inserts another file (which need not start with \code{\%!GLoBES}); includes
can be nested ten levels deep. Several experiments are separated by
\code{\#NEXT\#}, which resets all names, so that
\begin{aedl}
%!GLoBES
include "NOvA.glb"
#NEXT#
include "T2K.glb"
\end{aedl}
defines two experiments, loaded with \code{load\_experiments}.

\paragraph{Where files are found.} Included files, flux files and cross
section files are looked up in the current directory, in the directory of the
AEDL file that names them, and in the directories of \code{GLB\_PATH}.

\begin{cdiff}
C GLoBES looks for data files only in the current directory and in
\code{GLB\_PATH}, not next to the AEDL file, so that loading an experiment from
another directory fails to find its flux and cross section files.
\end{cdiff}
